\documentclass[11pt,a4paper]{article}
\usepackage[T1]{fontenc}
\usepackage[utf8]{inputenc}
\usepackage{lmodern}
\usepackage[margin=25mm]{geometry}
\usepackage{amsmath,amssymb,amsthm,mathtools,booktabs,array,longtable}
\usepackage{microtype,enumitem,xcolor,fancyhdr,xurl,hyperref,listings}
\hypersetup{colorlinks=true,linkcolor=blue!45!black,citecolor=blue!45!black,
 urlcolor=blue!45!black,pdftitle={Twelve Integer Variables through Interpolated Pell Certificates},
 pdfauthor={Research contribution for Vladimir Reshetnikov}}
\setlength{\headheight}{14pt}
\pagestyle{fancy}\fancyhf{}\fancyhead[L]{\small Interpolated Pell certificates}
\fancyhead[R]{\small ProveIt research contribution}\fancyfoot[C]{\thepage}
\setlength{\parskip}{4pt}\setlength{\parindent}{0pt}\emergencystretch=2em
\setlist{nosep,leftmargin=1.6em}
\newtheorem{theorem}{Theorem}[section]
\newtheorem{lemma}[theorem]{Lemma}
\newtheorem{proposition}[theorem]{Proposition}
\newtheorem{corollary}[theorem]{Corollary}
\theoremstyle{definition}\newtheorem{definition}[theorem]{Definition}
\theoremstyle{remark}\newtheorem{remark}[theorem]{Remark}
\newcommand{\Z}{\mathbb Z}\newcommand{\Q}{\mathbb Q}\newcommand{\N}{\mathbb N}\newcommand{\R}{\mathbb R}
\newcommand{\Con}{\operatorname{Con}}\newcommand{\F}{\mathcal F}
\newcommand{\Dm}{D_{\mathrm m}}\newcommand{\Df}{D_{\mathrm f}}
\newcommand{\Dt}{D_{\tau}}\newcommand{\Da}{D_{\mathrm a}}
\newcommand{\Dg}{D_{\mathrm g}}\newcommand{\Pair}{\operatorname{Pair}}
\newcommand{\NF}{\operatorname{NF}}
\numberwithin{equation}{section}
\lstset{basicstyle=\ttfamily\small,breaklines=true,columns=fullflexible,
 keepspaces=true,frame=single,rulecolor=black!20}
\title{\vspace{-8mm}\textbf{Twelve Integer Variables through\\Interpolated Pell Certificates}\\[3mm]
\large A monotone intersective merger, trace before separation,\\
and explicit Diophantine equations for syntactic consistency}
\author{Research contribution prepared for\\Vladimir Reshetnikov's ProveIt project}
\date{24 October 2026\\\small Reviewed repository commit:
\texttt{04a1d79222d55cabe4ba52897863fc33baa8c268}}
\begin{document}\maketitle
\begin{abstract}
We construct an ordinary polynomial in twelve unrestricted integer variables,
with exact total degree $4832$ and $248$ binary arithmetic operations, whose lack of integer solutions is equivalent
in ZFC to syntactic $\Con(\mathrm{ZFC})$, by composition with explicitly named
ordinary calibration theorems already used in ProveIt. The new reduction
shares a positivity witness with a strong Pell certificate. Full residue
interpolation of the auxiliary Pell coordinate then lets that same coordinate
enforce a divisibility condition and exclude a zero multiplier. A strictly
increasing intersective quintic,
$(x^2+31)(x^3+x+1)$, performs this merger while the main Pell square root is
still algebraic. Its absence of roots in real quadratic fields and the
strict positivity of its divided difference are essential.

A second construction uses thirteen unrestricted integers, exact degree
$1072$, and $190$ binary arithmetic operations, improving the inspected
thirteen-integer $1088$-degree/$195$-operation construction in both measures.
Its new ingredient is a normalized field trace that recovers a congruence
before square-root separation. Both polynomials have fully specified fixed
coefficients, literal division-free arithmetic graphs, and offline checks.
We make no claim of global minimality, a world variable-count record, or
historical priority for every constituent idea. The new proofs are ordinary
mathematical proofs; this package does not add an end-to-end Lean or encoded
object-language ZFC derivation.
\end{abstract}
\clearpage\tableofcontents\clearpage

\section{What is improved, and what is meant by an equation}

The task is to specify a polynomial $P\in\Z[x_1,\ldots,x_m]$, together with
its domain, for which
\begin{equation}
 \mathrm{ZFC}\vdash
 \bigl[\neg\exists\mathbf x\in\Z^m\;P(\mathbf x)=0\bigr]
 \longleftrightarrow\Con(\mathrm{ZFC}).
 \label{eq:target}
\end{equation}
The consistency sentence refers to the recursive finite-proof coding in
the inherited calibration. Proving this biconditional does not prove
either side of it. No coefficient is selected by deciding consistency,
waiting for a proof search, or choosing an unspecified enumeration index.

There is no single numerical meaning of ``simplest.'' A polynomial with
fewer variables may have much larger degree, a shorter circuit may use
large constants, and an expanded expression may be enormous despite a
small circuit. We preserve these distinctions. An arithmetic circuit here
has only integer leaves, the declared variable leaves, and binary
addition, subtraction and multiplication. Squaring costs a multiplication;
fixed-scalar multiplication is charged; reuse of a computed value is free.
Every intermediate letter in a formula below is a polynomial abbreviation,
not an additional quantified coordinate.

The comparison snapshot includes the requested Hilbert Tenth directory,
the Diophantine consistency collection, and all fifteen ZIP archives then
present in \texttt{docs/incoming}. Their hashes are recorded in the package.
The incoming root-fusion manuscript already contains a thirteen-integer
construction of degree $1088$ and $195$ operations. The active development
and other incoming proposals include ten-integer constructions of much
larger degree. Thus the twelve-variable result is a new tradeoff within
the inspected construction families, not a new universal variable bound.
Sun's recent ten-unknown representation theorem also rules out presenting
twelve unknowns as the current general record~\cite{sun2026}.

\begin{table}[ht]\centering
\begin{tabular}{lrrl}\toprule
Construction & Vars. & Degree & Arithmetic\\\midrule
Incoming five-root CRT norm &13&1088&195 operations\\
Trace-first companion, this package &13&1072&190 operations\\
Interpolated quintic merger, this package &12&4832&248 operations\\
Shorter twelve-variable companion &12&16000&223 operations\\
\bottomrule\end{tabular}
\caption{The directly relevant comparison. All domains are unrestricted
integers. The twelve-variable tradeoff spends degree to remove one more
coordinate. The machine-readable metrics report its full operation census.}
\end{table}

The main conceptual advance is not merely an arithmetic rewrite. The
former construction represents a strict positivity condition using three
square/triangular coordinates, one of which is absorbed into a norm. The
new construction uses a Pell positivity witness, shares its multiplier
with the strong certificate, and reuses the auxiliary coordinate as an
intersective residue. The proof must establish full residue coverage;
ordinary existence of Pell solutions would not suffice.

\section{Fixed data and the calibrated arithmetic interface}
\label{sec:interface}

Both constructions use the same coherent six-integer tuple:
\begin{equation}
 \beta=2^{1064921}-1,\quad K_0,\quad M_C,\quad M_F,
 \quad A_c=2^{1113},\quad B_c=2^{666104}.
 \label{eq:ports}
\end{equation}
The full finite recipe and all actual numeral bytes are included. In
particular $M_F=M_{F,\mathrm{native}}+\beta$ is already shifted. It must
not be shifted again. Appendix~\ref{app:coeff} prints the bounded recipe.
The ordered $943$-tile table and $430$ machine actions are supplied data,
so the recipe specifies actual integers rather than free parameters.

For the twelve-variable construction use precisely
\begin{equation}
 \boxed{j,C,z,h,r,v,w,T,c,n,t,s\in\Z.}
 \label{eq:variables12}
\end{equation}
Capital $C$ and lowercase $c$ are distinct coordinates. Define
\begin{align}
 J&=6j+1,&q&=\beta J+1,&W_0&=12w+4,\nonumber\\
 X&=qW_0,&Y&=(2v+1)q^3,&E&=XY,\nonumber\\
 a&=E+Y,&A&=a+2,&H&=4a+3,&\Delta&=A^2-1=a^2+H.
 \label{eq:front12}
\end{align}
Use the polynomial transport substitution
\begin{align}
 F&=(K_0+W_0)C-z(q-1),&\alpha&=q-F-C,\nonumber\\
 R&=\bigl(q(q-F)-C\bigr)(q^2-1)+(M_C+qM_F)J,\nonumber\\
 k&=R+1+hE,&\eta&=c-Yk,&\zeta&=k-\eta,\nonumber\\
 b&=X+ac,&\Dm&=1+\Delta c^2,\nonumber\\
 L_\tau&=XY^2k,&\Dt&=1+L_\tau(L_\tau+k).
 \label{eq:transport12}
\end{align}
The last factorization is the established identity
$XY^2(XY^2+1)k^2=L_\tau(L_\tau+k)$; its reuse is credited to the
existing Hilbert Tenth construction.

The two strict guards are combined by
\begin{align}
 U_c&=A_cF,&Z_c&=B_cC,\nonumber\\
 G&=U_cZ_c-(q-U_c-Z_c)^2,&P_\eta&=\eta\zeta,\nonumber\\
 O&=GP_\eta-(t^2-G-P_\eta+1)^2.
 \label{eq:pair12}
\end{align}
For integer $x,y$, there is an integer $t$ with
\begin{equation}
 \Pair(x,y,t):=xy-(t^2-x-y+1)^2>0
 \quad\Longleftrightarrow\quad x>0\text{ and }y>0.
 \label{eq:pairlemma}
\end{equation}
The forward implication follows because two negative factors make the
squared term larger than their product. For the converse put
$N=x+y-1$ and choose the nearest integer square $t^2$. If
$m=\lfloor\sqrt N\rfloor$, one of $m,m+1$ gives
$|t^2-N|\le m$, with strict squared error $<N\le xy$ unless the
error is zero. This proves the claim. The square root in this witness
selection is not an operation in the Diophantine polynomial.

The inherited signed Pell/native interface identifies finite ZFC
refutations with canonical solutions of the transport, first and main
Pell equations, the strict guards, and a strong/auxiliary rank certificate.
In the completeness direction it provides
\begin{equation}
 A>2\text{ even},\quad R>0,\quad R\equiv3\pmod4,\quad
 c=\psi_A(R)>0,\quad d=\chi_A(R)=b+H\gamma,
 \label{eq:canonical}
\end{equation}
with integral $\gamma$, together with $\tau^2=\Dt$ and $G,P_\eta>0$.
Here
\begin{equation}
 \chi_A(m)+\psi_A(m)\sqrt{A^2-1}
   =(A+\sqrt{A^2-1})^m.
 \label{eq:pellseq}
\end{equation}
The determinant guard uses canonical native witnesses; it is an
existence-preserving restriction, not a statement about every older witness.

Our strong square is already normalized: $f^2-\Delta I^2=1$ with
$I=src^2$. A direct rank argument in Section~\ref{sec:calibration}
therefore supplies the strong-index conclusion without extending an
odd-only source parametrization. The auxiliary congruences remain
$V\equiv-c\pmod f$ and $V\equiv-R\pmod c$. The new merger proves
$I\ne0$. No half-integer old coordinate is introduced.

\section{A monotone intersective quintic}
\label{sec:quintic}

Define
\begin{equation}
 \Pi(x)=(x^2+31)(x^3+x+1)
       =x^5+32x^3+x^2+31x+31.
 \label{eq:pi}
\end{equation}
The particular combination of properties below is what the construction
uses. The general technique of intersective divisibility mergers is
classical and is used in Sun's recent reduction; the inspected repository
also already uses the different quintic $(x^2+23)(x^3-x-1)$.

\begin{lemma}\label{lem:pi}
The polynomial $\Pi$ has a root modulo every positive integer, has no
root in any real extension of $\Q$ of degree at most two, and is strictly
increasing on $\R$.
\end{lemma}
\begin{proof}
The cubic $f=x^3+x+1$ has discriminant $-31$ and no rational root:
the only possible rational roots are $1,-1$, and neither is a root.
It is therefore irreducible over $\Q$. Its roots cannot lie in a field
of degree at most two. The quadratic $x^2+31$ has no real root.

For an odd prime $p\ne31$, either $f$ has a root in $\F_p$, or it is
an irreducible cubic. In the latter case Frobenius permutes its three
roots as a three-cycle. The Vandermonde product is fixed by that even
permutation, so its square, the discriminant $-31$, is a nonzero square
in $\F_p$. Thus $x^2+31$ has a simple root. In the former case a cubic
root is simple because the discriminant is nonzero modulo $p$.
At $p=31$, $x=3$ is a simple cubic root, since $f(3)=31$ and
$f'(3)=28$. These simple roots lift to every power of the corresponding
odd prime. At $2$, $-31\equiv1\pmod8$, and an odd square root lifts
successively to every $2$-power. The Chinese remainder theorem supplies
a root modulo any positive integer.

Finally
\[
 \Pi'(x)=5x^4+96x^2+2x+31
        =5x^4+95x^2+(x+1)^2+30>0.
\]
\end{proof}

For nonzero integer $H$ put
\begin{equation}
 \Theta(H,\xi,T)=H^5\Pi\left(\frac{\xi+HT}{H}\right).
 \label{eq:theta}
\end{equation}
This is a homogeneous integer polynomial of degree five in $H,\xi+HT$.
If $\xi$ is integral and $H^5\mid\Theta(H,\xi,T)$, then $H\mid\xi$.
Indeed, if $p$ has $\nu_p(\xi)<\nu_p(H)$, the monic leading term
$(\xi+HT)^5$ has strictly smaller $p$-valuation than every other term.
Its valuation is $5\nu_p(\xi)<5\nu_p(H)$, a contradiction.

The merger must operate before the main root $d=\sqrt{\Dm}$ is known to
be integral. Set $u=HT-b$. Reducing $d^2=\Dm$ gives
\begin{equation}
 \Theta(H,d-b,T)=\mathcal A+\mathcal B d,
 \label{eq:AB}
\end{equation}
where
\begin{align}
 \mathcal A={}&u^5+10u^3\Dm+5u\Dm^2
       +32H^2(u^3+3u\Dm)+H^3(u^2+\Dm)
       +31H^4u+31H^5,\nonumber\\
 \mathcal B={}&5u^4+10u^2\Dm+\Dm^2
       +32H^2(3u^2+\Dm)+2H^3u+31H^4.
 \label{eq:ABexplicit}
\end{align}
Since $\Dm>0$ and $H\ne0$, the coefficient $\mathcal B$ is positive.
One purely algebraic certificate is
\[
 \mathcal B=5u^4+10u^2\Dm+\Dm^2+32H^2\Dm
 +95H^2u^2+H^2(u+H)^2+30H^4>0.
\]
Also $\mathcal A\pm\mathcal B\sqrt{\Dm}\ne0$, by
Lemma~\ref{lem:pi}, even if $\Dm$ is not a square. This excludes a
zero multiplier without first assuming squarehood.

\section{The twelve-variable polynomial}
\label{sec:polynomial12}

Keep all definitions in Section~\ref{sec:interface} and define
\begin{align}
 I&=src^2,&\Df&=1+\Delta I^2,&Q&=\Delta^2I^2,\nonumber\\
 V&=\Df(c-R+4IT)-c,&
 \Da&=(Q-1)(QV^2-1),\nonumber\\
 \Dg&=(3O-4)s^2+1,&S&=IH^5,&Z&=Sn-\mathcal A.
 \label{eq:strong12}
\end{align}
The two uses of $T$ are intentional: it is both the auxiliary Pell
coordinate in $V$ and the intersective residue in $u=HT-b$.
Put $\ell_0=XY^2$. Use bounds on the root magnitudes, with one bound
for each stage:
\begin{align}
 K_T&=T^2+c^2+1,&B_0&=256K_T^2,\nonumber\\
 B_b&=1+c^2+R+c^2(s^2+r^2)(T^2+1),\nonumber\\
 B_{\rm a}&=Q(\Df B_b+c^2+1),&
 B_{\rm f}&=\Delta+I^2+3,&B_{\tau}&=2+\ell_0^2+k^2.
 \label{eq:bounds12}
\end{align}
The four separating weights are
\begin{align}
 W_1&=2B_0+1,&W_2&=W_1(2B_{\rm a}+1),\nonumber\\
 W_3&=W_2(2B_{\rm f}+1),&
 W_4&=W_3(2B_\tau+1).
 \label{eq:weight12}
\end{align}
They are odd, hence nonzero, on every integer assignment. Their positivity
will be proved from a vanishing norm factor before using them to separate
roots. Define the five integer polynomial radicands
\begin{align}
 T_0&=\mathcal B^2\Dm,&T_1&=S^2W_1^2\Da,\nonumber\\
 T_2&=S^2W_2^2\Df,&T_3&=S^2W_3^2\Dt,&T_4&=S^2W_4^2\Dg.
 \label{eq:radicands12}
\end{align}
For formal radicals let
\begin{equation}
 \F_{32}(z;t_0,t_1,t_2,t_3,t_4)
   =\prod_{\varepsilon\in\{-1,1\}^5}
       \left(z-\sum_{i=0}^4\varepsilon_i\sqrt{t_i}\right).
 \label{eq:norm32}
\end{equation}
The sign symmetries make this an element of
$\Z[z,t_0,\ldots,t_4]$. The requested equation is
\begin{equation}
 \boxed{P_{12}(j,C,z,h,r,v,w,T,c,n,t,s)
       :=\F_{32}(Z;T_0,T_1,T_2,T_3,T_4)=0.}
 \label{eq:main12}
\end{equation}
There is no division in this formula and no variable exponent. The
radical product defines a finite polynomial. The delivered literal graph
implements it using only $+,-,\times$; a small recurrence is supplied
in Section~\ref{sec:normcircuit}.

\section{Soundness of the twelve-variable reduction}
\label{sec:soundness12}

\begin{lemma}[Adjacent-coefficient descent]\label{lem:descent}
For integers $Q>1,V$ and rational $e$, the equation
$QV^2-(Q-1)e^2=1$ implies $e\in\Z$.
\end{lemma}
\begin{proof}
Write $e=a/d$ in lowest terms. Then $d^2\mid Q-1$. If $|V|=1$,
then $|e|=1$. Otherwise take positive absolute values and set
\[
 V'=(2Q-1)V-2(Q-1)e,\qquad e'=(2Q-1)e-2QV.
\]
The norm is unchanged; $V'$ is integral; and the reduced denominator of
$e'$ remains $d$, since $2Q-1\equiv1\pmod d$. The inequalities
$e>V$ and $(Q-1)e<QV$ give $-V<V'<V$. The norm excludes $V'=0$.
Descent in the positive integer $|V|$ reaches $1$ without changing the
denominator, which must consequently have been $1$.
\end{proof}

\begin{lemma}[Separation over a multiquadratic field]\label{lem:separation}
If nonzero complex square roots of rational numbers have strictly
superincreasing absolute values and their signed sum is rational, every
one of those roots is rational.
\end{lemma}
\begin{proof}
If a root is irrational, some automorphism of their multiquadratic
splitting field changes its sign. Subtracting the conjugate rational
relation yields a nonempty vanishing signed subsum. Its largest term
has absolute value greater than the sum of the absolute values of the
others, contradicting the triangle inequality. Square classes need not
be independent, and negative radicands are allowed.
\end{proof}

For every integer assignment to the front end, $q$ is even,
$q\equiv2\pmod3$, $|q|\ge16$, $|X|\ge64$ and $|Y|\ge4096$.
Thus $|a|\ge63\cdot4096$, $H\ne0$ and $\Delta>1$.
The radicands $\Dm,\Df,\Dt$ are positive: for the last use
$u(u+1)\ge0$ for the integer $u=XY^2$.

Suppose $P_{12}=0$. If $S=0$, the four radicands $T_1,\ldots,T_4$
are zero and the polynomial specializes to
\[
 (\mathcal A^2-\mathcal B^2\Dm)^{16}.
\]
Both factors $\mathcal A\pm\mathcal B\sqrt{\Dm}$ are nonzero by
Lemma~\ref{lem:pi}. This is a contradiction. Therefore
\begin{equation}
 S\ne0,\qquad I\ne0,\qquad s\ne0,\quad r\ne0,\quad c\ne0.
 \label{eq:nonzero12}
\end{equation}
In particular $Q>1$ and $\Df=1+\Delta s^2r^2c^4>|c|+1$.
Writing $b'=c-R+4IT\in\Z$, either $b'=0$ and $|V|=|c|\ge1$,
or $|V|\ge\Df-|c|>1$. Hence $\Da>0$.

We first recover positivity without assuming any individual root is
rational. The weights $W_i$ are odd, hence nonzero. All radicands
except $\Dg$ are positive. If $\Dg<0$, its radical would be the only
nonzero imaginary term in a vanishing factor, which is impossible.
Thus $\Dg\ge0$. Since $s\ne0$, its defining expression forces $O>0$:
if $O\le0$, then $\Dg\le1-4s^2<0$. The pairing lemma gives
$G>0$ and $\eta\zeta>0$.

We also need $R>0$ before bounding the roots. The determinant inequality
implies that $F,C,\alpha$ have the sign of $q$. To see this, $U_c,Z_c$
have the same sign and
$|U_c+Z_c|<2|q|$; since $A_c,B_c\ge2$, we get
$|F|+|C|<|q|$. If $q>0$, then
$q(q-F)-C=q\alpha+(q-1)C>0$, and the mask term in $R$ is positive.
If $q=-q_1<0$, write $C=-C_1$, $\alpha=-\alpha_1$ with
$C_1,\alpha_1>0$. Then
$q(q-F)-C=q_1(C_1+\alpha_1)+C_1>0$.
Also $J<0$ and $M_C+qM_F=M_C-q_1M_F<0$, because the actual
fixed tuple has $M_F>\beta>M_C>0$. Thus $R>0$ in this case as well.

The three bounds $B_{\rm a},B_{\rm f},B_\tau$ are now positive.
The inequality
\[
 4|srT|\le(s^2+r^2)(T^2+1)
\]
gives $|c-R+4IT|\le B_b$ and $|V|\le\Df B_b+c^2+1$.
Consequently
\begin{equation}
 \sqrt{\Da}<Q|V|\le B_{\rm a},\qquad
 \sqrt{\Df}\le B_{\rm f},\qquad
 \sqrt{\Dt}\le B_\tau.
 \label{eq:rootbounds12}
\end{equation}
For the last two, use
$\sqrt{1+\Delta I^2}\le1+(\Delta+I^2)/2$ and
$\sqrt{\ell_0(\ell_0+1)}\le|\ell_0|+1$, together with
$2|xy|\le x^2+y^2$.

The main radical admits a much smaller bound than its polynomial
numerator suggests. Put $v_0=u/H$ and $z_0=\sqrt{\Dm}/|H|$.
The uniform front-end size bounds imply
\[
 \left|\frac bH\right|<1+\frac{|c|}{3},\qquad
 z_0<\frac{|c|}{2}.
\]
Indeed $|a/H|<1/3$, $|X/H|<1$ and
$\sqrt{1+\Delta c^2}<(|a|+3)|c|$, while
$|H|\ge4|a|-3>2(|a|+3)$.
Thus $v_0^2\le3K_T$ and $z_0^2\le K_T$. Homogeneity of the quintic gives
\begin{align*}
 \frac{\mathcal B}{H^4}
 &=5v_0^4+10v_0^2z_0^2+z_0^4+96v_0^2+32z_0^2+2v_0+31\\
 &\le76K_T^2+323K_T+32\le431K_T^2.
\end{align*}
Since $|I|=|sr|c^2\ge c^2$,
\begin{equation}
 \frac{\mathcal B\sqrt{\Dm}}{|S|}
 =\frac{\mathcal B}{H^4}\frac{z_0}{|I|}
 <\frac{431}{2}K_T^2<256K_T^2=B_0.
 \label{eq:homogeneousbound}
\end{equation}
This uses the denominator already present in the congruence merger; it
adds no coordinate and no division to the polynomial.

Divide a vanishing norm factor by $S$. The five magnitudes, in order, are
\[
 \frac{\mathcal B\sqrt{\Dm}}{|S|},\quad
 W_1\sqrt{\Da},\quad W_2\sqrt{\Df},\quad
 W_3\sqrt{\Dt},\quad W_4\sqrt{\Dg}.
\]
The first is less than $W_1/2$. After adding the auxiliary term, the sum
of the first two is less than
$W_1(1/2+B_{\rm a})=W_2/2$.
The same induction bounds the sum of the first three by $W_3/2$
and the first four by $W_4/2$.
Every nonzero root of an integer radicand has magnitude at least $1$.
Hence these terms are strictly superincreasing, after a zero guard term
is omitted. No upper bound on the last guard radical is needed.
Lemma~\ref{lem:separation} makes all five roots rational. Their integer
radicands and $\mathcal B\ne0$ give signed integers
\[
 d^2=\Dm,\quad a_1^2=\Da,\quad g^2=\Dg,\quad
 \tau^2=\Dt,\quad f^2=\Df.
\]
By Lemma~\ref{lem:descent}, $e=a_1/(Q-1)$ is integral and
$QV^2-(Q-1)e^2=1$.

The integral factor relation now implies
\begin{equation}
 S\mid\mathcal A+\mathcal B d
      =\Theta(H,d-b,T).
 \label{eq:div12}
\end{equation}
The monic valuation argument following \eqref{eq:theta} gives $H\mid d-b$,
so $\gamma=(d-b)/H$ is an integer and the main congruence is restored.
The previously recovered guards $G>0$ and $\eta\zeta>0$ remain available.

Set $\rho=\Delta I$. Then
\[
 \rho^2=\Delta(f^2-1),\qquad c^2\mid\rho,
 \qquad V\equiv-c\pmod f,\qquad V\equiv-R\pmod c.
\]
For the last congruence use $c^2\mid I$ and
$\Df\equiv1\pmod c$. The direct normalized rank argument in Section~\ref{sec:calibration}
now recovers the signed native interface. Consequently a polynomial zero
gives a finite refutation under the inherited native calibration.

\section{Full residue interpolation and completeness}
\label{sec:interpolation}

The crucial converse allows $T$ to take every residue modulo a fixed
strong witness $I$. Choosing a separate $I$ after choosing each residue
would not prove the statement needed by the merger.

\begin{lemma}[Positivity by a Pell square]\label{lem:positivity}
For every integer $O$,
\[
 O>0\quad\Longleftrightarrow\quad
 \exists s\in\Z\setminus\{0\}\;[(3O-4)s^2+1\text{ is a square}].
\]
\end{lemma}
\begin{proof}
For $O\le0$ the expression is negative. At $O=1$, take $s=1$,
giving zero. For $O\ge2$, the positive integer $3O-4$ is $2$ modulo
$3$, hence nonsquare. The ordinary Pell equation supplies a nontrivial
solution. This is the positivity device used in Sun's relation-combining
argument~\cite{sun2026}.
\end{proof}

\begin{theorem}[Native auxiliary interpolation]\label{thm:interpolation}
Let $A>2$ be even, $\Delta=A^2-1$, $R>0$ with $R\equiv3\pmod4$,
and $c=\psi_A(R)$. Given any positive integer $N$, there is a positive
integer $I$ divisible by $Nc^2$ and an integer $f>0$ satisfying
$f^2-\Delta I^2=1$, such that for every residue $a_0\pmod I$ there
is an integer $T$ for which
\[
 V=f^2(c-R+4IT)-c,\qquad
 (\Delta I)^2V^2-((\Delta I)^2-1)e^2=1
\]
has an integer $e$. The chosen $T$ can be made arbitrarily large and
satisfies $T\equiv a_0\pmod I$.
\end{theorem}
\begin{proof}
Finite order of the Pell matrix modulo $4Nc^2$ gives an index $m>0$
with $I=\psi_A(m)$ divisible by $4Nc^2$. Put $f=\chi_A(m)$.
In particular $m$ is even, $f$ is odd and $\gcd(f,I)=1$.
The standard Lucas valuation formulas give
\begin{equation}
 p\mid I\quad\Longrightarrow\quad\nu_p(m)\le\nu_p(I).
 \label{eq:valuation}
\end{equation}
For odd primes these are the rank-of-apparition valuation formulas;
at $2$, since $A$ is even,
$\nu_2(\psi_A(m))=\nu_2(A)+\nu_2(m)$. These are exactly the
specializations of the formulas used in the full-residue argument of
Sun~\cite{sun2026,sanna}. It follows that
\[
 \gcd(4mf,4I^2)\mid4I.
\]
The factor $f$ causes no obstruction because it is coprime to $I$.

Choose an arbitrarily large positive $\ell$ by the Chinese remainder
theorem with
\begin{equation}
 \ell\equiv R\pmod{4mf},\qquad
 \ell\equiv R-4Ia_0\pmod{4I^2}.
 \label{eq:ellcrt}
\end{equation}
In particular $\ell\equiv3\pmod4$. The period $4mf$, rather than
merely $4m$, is necessary modulo $f^2$. Indeed for
$\alpha=A+\sqrt\Delta$ we have
$\alpha^{2m}=2f\alpha^m-1$. Squaring and then raising to the
$f$th power gives $\alpha^{4mf}\equiv1\pmod{f^2}$ in
$\Z[\sqrt\Delta]$.

Set $D=\Delta I$ and
\[
 V_0=\chi_D(\ell)/D,\qquad e=\psi_D(\ell).
\]
For odd $\ell$, $\chi_D(\ell)$ is divisible by $D$, and the Pell
identity gives $D^2V_0^2-(D^2-1)e^2=1$.
The odd Chebyshev identity, using $D^2\equiv-\Delta\pmod{f^2}$,
gives
\[
 V_0\equiv-\psi_A(\ell)\equiv-c\pmod{f^2}.
\]
For $\ell\equiv3\pmod4$, the same odd polynomial has constant
term $-\ell$ and every nonconstant term is divisible by $4D^2$;
hence
\[
 V_0\equiv-\ell\pmod{4I^2}.
\]
Since $A$ is even, the Pell recurrence modulo $4$ gives
$c=\psi_A(R)\equiv R\pmod4$. Also
$f^2=1+\Delta I^2$, so
\[
 V_0+c-f^2(c-R)\equiv R-\ell\equiv4Ia_0\pmod{4I^2}.
\]
This numerator is divisible by $f^2$ and by $4I$. As $f$ is odd
and coprime to $I$, define the integer
\[
 T=\frac{V_0+c-f^2(c-R)}{4If^2}.
\]
Dividing the preceding congruence by $4I$, and using
$f^2\equiv1\pmod I$, gives $T\equiv a_0\pmod I$.
As $\ell$ tends to infinity through the CRT progression, $V_0$ and
then $T$ tend to infinity. This proves full residue coverage for the
single chosen pair $(I,f)$.
\end{proof}

\begin{theorem}[Completeness of the twelve-variable equation]
Every canonical native refutation witness extends to an integer zero
of $P_{12}$.
\end{theorem}
\begin{proof}
Retain its front-end coordinates and its roots $d,\tau$, and first
choose $t$ with $O>0$. Lemma~\ref{lem:positivity} gives $s>0$ and
an integer $g$ with $g^2=(3O-4)s^2+1$.
Apply Theorem~\ref{thm:interpolation} with $N=s$. It gives a fixed
$I$ divisible by $sc^2$; put $r=I/(sc^2)$. Write
$\gamma=(d-b)/H\in\Z$.

By Lemma~\ref{lem:pi}, choose $a_0\pmod I$ with
$\Pi(\gamma+a_0)\equiv0\pmod I$. Use the interpolation theorem
to choose the same auxiliary $T$ in that residue. Then
\[
 \Theta(H,d-b,T)=H^5\Pi(\gamma+T)
\]
is divisible by $S=IH^5$. The auxiliary root gives
$a_1=(Q-1)e$ with $a_1^2=\Da$.
All five radicals in \eqref{eq:radicands12} now have integral roots.
Choose
\[
 n=\frac{\mathcal A+\mathcal B d}{S}
       +W_1a_1+W_2f+W_3\tau+W_4g.
\]
Every quantity on the right has already been chosen, and the quotient
is integral. It makes one norm factor vanish. There is no circular
selection of $I,T,n$, and no extra divisibility promise on the canonical
native data.
\end{proof}

\section{Exact degree and the literal norm circuit}
\label{sec:normcircuit}

For three radicals use the eight-factor norm
\begin{equation}
 \F_8(x;A,B,C)
 =\bigl((x^2-A-B-C)^2-4(AB+AC+BC)\bigr)^2-64ABCx^2.
 \label{eq:norm8}
\end{equation}
Compute $\F_8(x-u-v;A,B,C)$ in the rank-four algebra
with $u^2=D$, $v^2=E$, using the basis $1,u,v,uv$.
If the resulting vector is $(f_0,f_1,f_2,f_3)$, set
\begin{align*}
 g_0&=f_0^2+Ef_2^2-D(f_1^2+Ef_3^2),\\
 g_1&=2(f_0f_2-Df_1f_3),\\
 \F_{32}(x;A,B,C,D,E)&=g_0^2-Eg_1^2.
\end{align*}
Every operation is polynomial. The archive retains the small inherited
norm implementation and checks it against an independent formal
conjugate product. The twelve-variable graph specializes the cleared
norm routine at denominator $1$; no denominator leaf remains.

The important front-end degrees of $P_{12}$ are
\[
\begin{array}{c|rrrrrrrrrr}
\text{register}&a&H&\Delta&I&\Df&Q&V&\Da&O&\Dg\\\hline
\deg&6&6&12&4&20&32&25&114&44&46.
\end{array}
\]
Also $\deg\mathcal A=35$, $\deg\mathcal B=28$,
$\deg S=34$ and $\deg Z=35$.
The bound degrees are $4,58,12,20$ for
$B_0,B_{\rm a},B_{\rm f},B_\tau$, respectively.
Thus $\deg(W_1,W_2,W_3,W_4)=(4,62,74,94)$.
The five radicands have respective degrees
\[
70,\quad190,\quad212,\quad250,\quad302.
\]
The last is uniquely largest, and $\deg Z=35<151$.
The unique highest contribution in the norm is therefore $T_4^{16}$,
of degree $4832$.

For a specialization-proof nonzero coefficient put
\[
 q_0=6\beta j,\qquad Y_0=2vq_0^3,\qquad
 a_0=24wvq_0^4,\qquad I_0=src^2.
\]
The leading form of the guard is
$(\Dg)_{46}=-3s^2h^4a_0^4Y_0^4$.
Together with the weight definitions this gives
\[
 (P_{12})_{4832}
 =2^{704}3^{16}a_0^{544}Y_0^{128}I_0^{160}
   c^{64}s^{32}h^{64}(T^2+c^2)^{64}(s^2+r^2)^{32}T^{64}.
\]
Selecting $T^{128}$ from the first binomial and $r^{64}$ from the
second gives
\begin{equation}
 [j^{2560}w^{544}v^{672}h^{64}s^{192}r^{224}c^{384}T^{192}]P_{12}
   =2^{5024}3^{3120}\beta^{2560}>0.
 \label{eq:degree12}
\end{equation}
The variable exponents sum to $4832$. This is an exact degree certificate
on the actual fixed coefficient tuple.

\subsection{A shorter circuit at degree 16000}

There is a useful operation--degree tradeoff. Use the same front end,
$S,Z,\mathcal A,\mathcal B$, but set
\[
 B_*^{\rm short}=1+\mathcal B^2\Dm+\Dt+\Df+\Da+\Dg^2,
 \qquad \mathcal W^{\rm short}=2B_*^{\rm short}+1.
\]
and replace the five radicands by
\[
 \mathcal B^2\Dm,\quad
 S^2(\mathcal W^{\rm short})^2\Da,\quad
 S^2(\mathcal W^{\rm short})^4\Dg,\quad
 S^2(\mathcal W^{\rm short})^6\Dt,\quad
 S^2(\mathcal W^{\rm short})^8\Df.
\]
Use these in the same $\F_{32}$ norm centered at $Z$.
The same proof works; even before the guard has been recovered,
this bounds all five unweighted absolute magnitudes once $S\ne0$.
Its literal graph costs $223$ binary operations. Here
$\deg\mathcal W^{\rm short}=114$, and the five radicand degrees are
$70,410,570,786,1000$. The exact degree is $16000$.
Writing
\[
 R_0=-q_0^3(12wC-q_0z),\qquad b'_0=4I_0T-R_0,
\]
its highest homogeneous part is
\[
 2^{448}a_0^{1728}I_0^{1088}(b'_0)^{256}.
\]
In particular the coefficient of
$j^{6912}w^{1728}v^{1728}s^{1344}r^{1344}c^{2688}T^{256}$
is $2^{13056}3^{8640}\beta^{6912}>0$.
Both versions have exactly the same twelve integer coordinates, the same
six coefficient numerals, and the same solvability equivalence.
The archive supplies both literal graphs.

\section{Thirteen integers: trace before separation}
\label{sec:trace13}

The companion equation has thirteen unrestricted integer coordinates,
exact degree $1072$, and $190$ live binary operations. It improves the
inspected CRT companion of \cite{fusion}, which has degree $1088$ and
$195$ operations, while preserving its complete fixed coefficient tuple.
The key step recovers a divisibility condition from a norm zero before
requiring individual radicals to be rational. Consequently the original
large-scale separation bounds become available earlier, and two scaling
devices can be removed.

\subsection{Complete specification}

Use the six constants \eqref{eq:ports} and exactly the coordinates
\begin{equation}
 \boxed{j,C,z,h,r,v,w,T,c,n,t,p_2,p_3\in\Z.}
 \label{eq:trace13variables}
\end{equation}
The common front end is recomputed with $Y=qv$:
\begin{align}
 J&=6j+1,&q&=\beta J+1,&W_0&=12w+4,\nonumber\\
 X&=qW_0,&Y&=qv,&E&=XY,\nonumber\\
 a&=E+Y,&H&=4a+3,&\Delta&=a^2+H,\nonumber\\
 F&=(K_0+W_0)C-z(q-1),&\alpha&=q-F-C,\nonumber\\
 R&=\bigl(q(q-F)-C\bigr)(q^2-1)+(M_C+qM_F)J,\nonumber\\
 k&=R+1+hE,&\eta&=c-Yk,&\zeta&=k-\eta.
 \label{eq:trace13front}
\end{align}
Let $G,O$ have their exact definitions in \eqref{eq:pair12}, using these
new values of $F,q,\eta,\zeta$. Define
\begin{align}
 I&=(2r+1)c^2,&\Df&=1+\Delta I^2,&Q&=\Delta^2I^2,\nonumber\\
 \delta&=Q-1,&V&=\Df(cT-R)-c,&C_y&=QV^2-1,\nonumber\\
 L_\tau&=XY^2k,&\Dt&=1+L_\tau(L_\tau+k),\nonumber\\
 \Dm&=1+\Delta c^2,&b&=X+ac,\nonumber\\
 A_g&=O-1-p_2^2-p_3(p_3+1),\nonumber\\
 M&=2q^2,&U&=MH,&
 Z&=U(n+V)+Mb+H(v-q^2).
 \label{eq:trace13registers}
\end{align}
In particular this section's $I,V,A_g$ differ from the twelve-variable
definitions. There is no multiplier $v^2$ in $A_g$ and no factor
$4H+1$ on the main radical. Using the polynomial \eqref{eq:norm32}, put
\begin{equation}
 \boxed{P_{13}^{\rm tr}=
 \delta^{16}\F_{32}\left(
 Z;M^2\Dm,U^2\Df,U^2\Dt,U^2A_g,\frac{U^2C_y}{\delta}
 \right)=0.}
 \label{eq:trace13poly}
\end{equation}
The degree in the final radicand of $\F_{32}$ is at most sixteen, so
the displayed denominator clears symbolically. The delivered arithmetic
graph contains only integer leaves and binary $+,-,\times$; neither
division nor radical extraction is an evaluator operation.

\subsection{A congruence from the rational part}

\begin{lemma}[Rational-part divisibility]\label{lem:trace13rational}
Let $a_i\in\Q$, possibly zero or negative, and choose roots
$\alpha_i^2=a_i$. Suppose every rational $\alpha_i$ is integral.
If $Z,L_1,\ldots,L_m\in\Z$ and
$Z=\sum_iL_i\alpha_i$, then $\gcd(L_1,\ldots,L_m)$ divides $Z$.
\end{lemma}
\begin{proof}
In the finite multiquadratic splitting field
$K=\Q(\alpha_1,\ldots,\alpha_m)$ apply
$\pi(x)=\operatorname{Tr}_{K/\Q}(x)/[K:\Q]$.
A rational root is fixed and contributes its integer value. An
irrational root has minimal polynomial $x^2-a_i$, of trace zero;
trace transitivity therefore gives $\pi(\alpha_i)=0$. Thus
$Z=\sum_iL_i\pi(\alpha_i)$ is an integer linear combination of the
$L_i$. Imaginary roots also have rational part zero. Independence of
square classes and separation of magnitudes are unnecessary.
\end{proof}

Here the four radicands $\Dm,\Df,\Dt,A_g$ are integers. For the
remaining rational radicand $D_y=C_y/\delta$, a rational root is
integral by Lemma~\ref{lem:descent} when $Q>1$; when $Q=0$ it equals
$\pm1$. The following checks establish that these cases exhaust every
integer assignment, before imposing the equation.

Always $q$ is even, $q\equiv2\pmod3$, $|q|\ge16$, and
$|X|\ge64$. At $v=0$ we have $Y=a=0$, $H=3$, and $\Delta=3$.
Otherwise $|Y|\ge16$ and $|a|\ge63\cdot16$. Hence in both cases
$H\ne0$ is odd and $\Delta\ge3$ is positive and odd. The radicands
$\Df,\Dm$ are at least one. Writing $u=XY^2\in\Z$ gives
$\Dt=1+u(u+1)k^2\ge1$.

At $c=0$, $Q=0,\delta=-1,C_y=-1$, and $D_y=1$. At $c\ne0$,
$2r+1\ne0$ gives $Q\ge9$ and $\Df>|c|+1$. For the integer
$b_1=cT-R$, either $b_1=0$ and $|V|=|c|\ge1$, or
$|V|\ge\Df-|c|>1$. Thus
\begin{equation}
 \delta\ne0,\qquad
 D_y=V^2+\frac{V^2-1}{Q-1}\ge1\quad(c\ne0).
 \label{eq:trace13realness}
\end{equation}
All four non-guard radicands are positive, including on the exceptional
$c=0$ and $v=0$ slices.

At a zero of \eqref{eq:trace13poly}, one conjugate factor vanishes.
Absorb its signs into roots $d,f,\tau,g,y$ to obtain
\begin{equation}
 Z=Md+U(f+\tau+g+y),\qquad
 (d^2,f^2,\tau^2,g^2,y^2)=(\Dm,\Df,\Dt,A_g,D_y).
 \label{eq:trace13factor}
\end{equation}
Lemma~\ref{lem:trace13rational} gives $M\mid Z$ immediately. This
does not assume that all five roots are rational or that $A_g\ge0$.
Since
\[
 Z=M\bigl(H(n+V)+b\bigr)+H(v-q^2),
\]
we have $M\mid H(v-q^2)$. But $Y=qv$ implies
$H\equiv3\pmod q$, and $3\nmid q$ with $H$ odd gives
\begin{equation}
 \gcd(H,M)=1,\qquad
 v=q^2(2v_0+1),\qquad
 Y=(2v_0+1)q^3,\quad
 v_0=\frac{v-q^2}{2q^2}\in\Z.
 \label{eq:trace13scale}
\end{equation}
Thus $v,Y\ne0$ and $|Y|\ge|q|^3\ge4096$ are established before
squarehood or Pell-index recovery. This early recovery is the new step.

\subsection{Unscaled guard, separation, and restoration}

If $A_g<0$, its radical in \eqref{eq:trace13factor} would be the only
nonzero imaginary term; therefore $A_g\ge0$. Consequently
$O=A_g+1+p_2^2+p_3(p_3+1)>0$, and \eqref{eq:pairlemma} gives
$G>0$ and $\eta\zeta>0$. This also excludes $c=0$, since there
$\eta\zeta=-Y(Y+1)k^2\le0$.

Put $x=|X|\ge64$, $m=|Y|\ge4096$, $K=|k|$, and $C_0=|c|$.
The strip guard gives $K\ge2$ and
\begin{equation}
 K(m-1)<C_0<K(m+1).
 \label{eq:trace13strip}
\end{equation}
Denote positive roots by $d_+,f_+,\tau_+,y_+$, and let
$g_+=\sqrt{A_g}\ge0$. With $a_*=|a|\ge(x-1)m$, the bounds
$\sqrt\Delta>a_*-3>xm/2$ yield
\begin{gather}
 f_+>\frac{xm}{2}C_0^2>3C_0^2,\qquad
 K(xm^2-1)<\tau_+<K(xm^2+1),\nonumber\\
 \tau_+>2C_0,\qquad f_+>2\tau_+,\qquad
 Q>\Df>\Dt>1,\quad Q>C_0^2,\nonumber\\
 \frac{C_0}{5}<\frac{d_+}{|H|}<\frac{C_0}{2}.
 \label{eq:trace13bounds}
\end{gather}
For example, $f_+>2\tau_+$ follows from
$mK(m-1)^2>4(m^2+1/x)$. The last comparison uses
$a_*-3<\sqrt\Delta<a_*+3$ and
$4a_*-3\le|H|\le4a_*+3$.
The determinant identity
\[
 12G=4q^2-\bigl(3(U_c+Z_c)-4q\bigr)^2-3(U_c-Z_c)^2
\]
gives $G<q^2$. Since $\eta\zeta\le k^2/4$, the unscaled guard gives
\begin{equation}
 g_+<\frac{|q|K}{2}<\frac{C_0}{10}.
 \label{eq:trace13guardbound}
\end{equation}
Here the final inequality uses the already recovered $m\ge|q|^3$.

If $b_1=cT-R=0$, then $V=-c$ and
$C_0<y_+<C_0+1$. Otherwise
$|V|\ge f_+^2-C_0>2f_+$ and $y_+>|V|>2f_+$.
After omitting a zero $g_+$, the following orders are strictly
superincreasing:
\begin{equation}
 \begin{array}{ll}
 b_1=0:&g_+,\ d_+/|H|,\ y_+,\ \tau_+,\ f_+,\\
 b_1\ne0:&g_+,\ d_+/|H|,\ \tau_+,\ f_+,\ y_+.
 \end{array}
 \label{eq:trace13orders}
\end{equation}
The two smallest sum to less than $0.6C_0$; in the first case the
three smallest sum to less than $1.6C_0+1<2C_0<\tau_+$.
The other comparisons follow from \eqref{eq:trace13bounds}.
Divide \eqref{eq:trace13factor} by $U$. Its main-root coefficient is
$M/U=1/H$, so Lemma~\ref{lem:separation} applies to precisely these
magnitudes. All roots are rational; the integer radicands make
$d,f,\tau,g$ integral, and Lemma~\ref{lem:descent} makes $y$ integral.

Reducing the integral relation \eqref{eq:trace13factor} modulo $H$
and using \eqref{eq:trace13scale} recovers
$\gamma=(d-b)/H\in\Z$. Thus we have the exact normalized kernel
\begin{equation}
 \begin{gathered}
 f^2=\Df,\quad\tau^2=\Dt,\quad
 QV^2-(Q-1)y^2=1,\\
 (b+H\gamma)^2=\Dm,\quad A_g=g^2,\quad
 Y=(2v_0+1)q^3.
 \end{gathered}
 \label{eq:trace13kernel}
\end{equation}
It is the incoming unscaled normalized integer kernel of \cite{fusion},
with its old scale coordinate $v_0$ and its guard square coordinate
replaced by $g$. Its transport substitution and entire fixed tuple
are unchanged.

Conversely, take any solution of that kernel, writing its guard root as
$p_1$ and its main root as $d=b+H\gamma$. Set
\begin{equation}
 v=q^2(2v_0+1),\qquad
 n=\gamma+f+\tau+p_1+y-V-v_0.
 \label{eq:trace13inverse}
\end{equation}
Every other supplied coordinate is retained. These integer choices give
\eqref{eq:trace13factor}, so one norm factor vanishes. The two directions
prove an existence equivalence with the unchanged kernel. Its ordinary
native/K215 interpretation is the inherited dependency described in
Section~\ref{sec:calibration}; it is not supplied by a finite norm test.

\subsection{Exact degree and complete arithmetic cost}

Removing $v^2$ from the old guard saves two gates; removing its optimized
main scale $4U+M$ saves two more. The established factorization
$\Dt=1+L_\tau(L_\tau+k)$ saves one further gate. The complete new
count is $118$ multiplications, $40$ additions and $32$ subtractions.
All thirteen variables, all six ports, and every one of the $190$ gates
are live. This includes the whole cleared norm.

For the degree calculation write
\[
 d=M^2\Dm,\quad e=U^2A_g,\quad A'=U^2\Df,
 \quad B'=U^2\Dt,\quad C'=U^2C_y,
\]
and $s_0=Un+Mb+H(v-q^2)$, so $Z=UV+s_0$. Their degrees are
$14,40,26,34,72$, while $\deg U=6$, $\deg V=19$,
$\deg\delta=22$, and $\deg s_0=7$. The first scalar component of the
biquadratic cleared three-root computation is exactly
\[
 J_0=\delta(s_0^2+2UVs_0+d+e-A'-B')-U^2(V^2-1).
\]
Giving $\sqrt d,\sqrt e$ weights $7,20$, exact sparse expansion gives
the successive four-component weighted degrees
\begin{equation}
 (62,54,67,49),\quad(134,116,129,121),\quad
 (268,250,263,255).
 \label{eq:trace13componentdegrees}
\end{equation}
The unique top scalar is $16\delta^4U^4V^4e^2$ of degree $268$.
The two remaining quadratic norms square it twice, giving degree $1072$.

For explicit attainment on the actual coefficient slice put
\[
 q_0=6\beta j,\quad a_0=12wvq_0^2,\quad i_0=2r,
 \quad R_0=q_0^4z-12q_0^3wC,\quad
 \eta_0=-q_0v(R_0+ha_0).
\]
The full leading form is
\begin{equation}
 (P_{13}^{\rm tr})_{1072}
 =2^{112}q_0^{64}a_0^{128}i_0^{64}c^{128}R_0^{16}\eta_0^{32}.
 \label{eq:trace13top}
\end{equation}
In particular,
\begin{equation}
 [j^{480}w^{160}v^{192}h^{32}r^{64}c^{128}z^{16}]
 P_{13}^{\rm tr}=2^{976}3^{640}\beta^{480}>0.
 \label{eq:trace13coefficient}
\end{equation}
The exponent sum is $1072$. The $h^{32}$ and then $z^{16}$ selections
are unique, so relations among the other five constants cannot cancel
this coefficient. The degree certificate checks these identities exactly;
modular evaluations separately check transcription of the complete graph.


\section{Consistency calibration and proof boundaries}
\label{sec:calibration}

\begin{theorem}[Ordinary consistency equivalence]
For the fixed tuple in Appendix~\ref{app:coeff}, both $P_{12}$ and
$P_{13}^{\mathrm{tr}}$ satisfy \eqref{eq:target}, with respectively
twelve and thirteen unrestricted integer coordinates.
\end{theorem}
\begin{proof}
For $P_{12}$, Section~\ref{sec:soundness12} restores the main and first
Pell roots, strict determinant and strip guards, a nonzero integer $I$,
and both auxiliary congruences. Write $B_0=|a+2|\ge2$ and
$c_0=|c|$. The inherited signed first/main rank bootstrap gives
$c_0=\psi_{B_0}(p)$ with $p\ge25$~\cite{signed,guards}.
The normalized strong equation $f^2-\Delta I^2=1$ directly gives
$|I|=\psi_{B_0}(m)$ and $|f|=\chi_{B_0}(m)$ for some $m>0$.
Since $c_0^2\mid|I|$, strong divisibility gives $p\mid m$.
Write $m=pk$, and put $C_0=\chi_{B_0}(p)$. Binomial expansion yields
\[
 \frac{\psi_{B_0}(pk)}{c_0}\equiv kC_0^{k-1}\pmod{c_0}.
\]
The left side is divisible by $c_0$, whereas $\gcd(C_0,c_0)=1$.
Thus $c_0\mid k$, and in particular $c_0\mid m$ and $m>2p$.
These are precisely the strong-index inputs for the inherited signed
auxiliary step-down. No generalized odd-multiplier lemma is required.

For that step-down put $\rho=\Delta I$. Then
\[
 \rho^2\equiv-\Delta\pmod{|f|},\qquad
 \rho^2\equiv0\pmod{c_0}.
\]
The exposed quotient equations also have integral polynomial witnesses:
\begin{align*}
 o&=f(c-R+4IT),\\
 j_{\rm aux}
   &=\Delta s^2r^2c^3(c-R)+4src\,\Df T,\\
 V&=fo-c=cj_{\rm aux}-R.
\end{align*}
The auxiliary Pell index is odd because
$\chi_{|\rho|}(2k)\equiv(-1)^k\pmod{|\rho|}$, independently of
the parity of $\rho$. The signed rank and ratio argument now gives
$p=R$, $2n_{\rm first}=R+1$, and $X=\operatorname{sign}(c)2^R$.
Since $X\equiv2\pmod3$, it selects $c>0$ and $X=2^R$.
The first-index sign identity gives $k>0$, hence $Y>0$.
Finally the Mersenne-radix sign lemma excludes negative $q$.
The post-power native decoder therefore applies directly
\cite{signed,guards}. These retained proofs cover the signed stages
before any positive-domain theorem is invoked.

Conversely, the inherited canonical native construction supplies
\eqref{eq:canonical}, including its even $A$ and $R\equiv3\pmod4$.
The interpolation theorem constructs a new strong and auxiliary witness
without altering those native coordinates. The completeness proof then
constructs the polynomial zero.

The trace companion has a direct two-way restoration to the incoming
normalized integer kernel, as proved in its section. The inherited
native-window/compiler theorem identifies the kernel's existence with
blank halting of the fixed K215 machine, and the retained ordinary K215
calibration identifies that halting with existence of a finite refutation
of the specified ZFC calculus. These are mathematical theorems used as
lemmas, with their exact source versions and explanatory proofs retained.
All new arguments use finite algebra, induction, elementary number theory,
Pell recurrences, finite fields and finite coding, available in ZFC.
Compose the existence equivalences and negate them.
\end{proof}

Three assertions must be distinguished. First, the new reductions have
the ordinary unbounded proofs supplied here. Second, the consistency
interpretation imports the named ordinary native and syntactic-calibration
theorems; this report does not reprove the entire compiler from its local
tile list. Third, executable audits certify finite identities, numeral
bindings, degrees and graph transcription. Those audits are not proofs of
the unbounded compiler theorem. The repository separately records that its
complete native Lean calibration and encoded object-ZFC derivation remain
unfinished~\cite{active,boundary}. This package does not change that status.

The inherited native/compiler and K215 texts supply ordinary proofs, not
conjectural placeholders. The dependency ledger identifies the exact lemmas
being composed. A future proof-assistant formalization must cover that
composition as well as the final polynomial algebra.

\section{Reproduction, attribution and further research}

The archive contains the three literal equations, builders, exact algebra
checks, separate mathematical audits, a manifest, all six fixed numeral
payloads, and their complete finite reconstruction. The retained
\texttt{inherited/root\_fusion} directory is an unchanged copy of the
inspected incoming source package. Its licenses, provenance and proof
boundaries remain attached. No existing repository files were overwritten
or published by this delivery.

The top-level verification entry point regenerates the graphs in memory,
checks the delivered equations and coefficients, runs independent
transcription diagnostics and negative controls, and writes an explicit
receipt. The precise reproduction commands and dependency requirements
are in \texttt{README.md}. Python optimization must not remove acceptance
checks. No verifier searches for a consistency witness, simulates an
unbounded machine execution, or treats a bounded lack of zeros as evidence
of consistency. The mathematical reviews are internal independent reviews
of this research session, not external human peer review.

\subsection{Further research questions}
\begin{enumerate}
\item \textbf{Reduce the separating weights further.} Bounding square-root
magnitudes instead of radicands already reduces the twelve-variable degree
from $16000$ to $4832$, using homogeneity and distinct weights. Further
source-specific bounds may improve the tradeoff again.
Any such change must retain the order of the proof: exclude $S=0$,
recover positivity, and then separate arbitrary integer assignments.
\item \textbf{Fuse square conditions through proved coprimality.} The
strong and auxiliary conditions might be combined into a single square
if the multiplier is chosen to absorb precisely the obstructing prime
factors. Any saving must include the resulting degree increase and must
preserve full residue interpolation of the shared $T$.
\item \textbf{Use trace before separation elsewhere.} The trace lemma
extracts congruences even when irrational radicals cancel. Search for
other redundant guard multipliers and scaling weights in the existing
nine-positive and ten-integer circuits. Rational radicands require a
separate integrality theorem; monicity cannot be assumed.
\item \textbf{Find smaller real-quadratic-safe mergers.} The monotone
quintic supplies intersectivity, real-quadratic nonvanishing and a
nonzero odd coefficient simultaneously. Can another polynomial reduce
the degree or number of operations while retaining all three properties?
An ordinary rational-root-free polynomial is insufficient here.
\item \textbf{Remove another witness through a proved shared family.}
Setting $I=sc^2$ would erase $r$, but would require simultaneous Pell
solvability with a prescribed positivity witness. Neither this report nor
the finite checks establish that assertion. A theorem furnishing such a
family would be a substantive further variable reduction.
\item \textbf{Complete the native formal certificate.} The most valuable
formalization target is the shared source-to-refutation calibration and
its internal-ZFC export. Certifying only the small final norm identities
does not close that semantic boundary.
\item \textbf{Compare alternative source encodings on equal terms.}
The direct numerical-calculus quartics already in the incoming collection
avoid K215 and Jones, at a cost in variables. The norm and interpolation
lemmas here are not intrinsically tied to K215; applying them to an
alternative calibrated source requires an actual complete source packet
and new counts, not an unspecified universal coefficient.
\end{enumerate}

\appendix
\section{Complete bounded coefficient recipe}\label{app:coeff}

The following is the full source-34 finite arithmetic recipe~\cite{coeff},
reproduced from the incoming root-fusion report with its licenses retained.
All ranges are bounded constants. The ordered \(943\) tile records
\((x_{i0},x_{i1},x_{i2},x_{i3})\), in west, north, east, south order,
are supplied in
\path{inherited/root_fusion/coefficients/recipe/wang943_tiles.json}.
Their ordering is part of the definition: selector \(0\) is Start
and selector \(1\) is Forbidden.
The supplied \texttt{k215.tm} contains all \(430\) working actions
from which the source builder reconstructs the whole ordered list.
Its SHA-256 is
\begin{center}\small\texttt{
fdbcef0825886035c3091c8537a69535bc96f9e265f0ae45863543d40ae70090}.
\end{center}
Thus this finite table is delivered data, not an unspecified machine
or a free coefficient parameter.

Put
\[
b=121,\quad L_0=8801,\quad d=bL_0=1064921,\quad
\rho=2^{121},\quad B_0=2^{1064921}.
\]
For \(0\le s<4\), \(0\le j<10\), define
\[
e_s(j)=o_s+2j,\qquad(o_0,o_1,o_2,o_3)=(1020,1000,980,960),
\]
and \(w_j=2^{10+2j}\) for \(0\le j<44\). Let
\[
\begin{split}
\mathcal E={}&\{0,\ldots,942\}
 \cup\{e_s(j):0\le s<4,\ 0\le j<10\}\\
 &\cup\{1063,1064,1095,1096,2195,3296\}.
\end{split}
\]
It has \(989\) elements and maximum \(3296\).
For each tile \(0\le i<943\), set
\[
c_i=1+\sum_{s=0}^3\sum_{j=0}^9
       \operatorname{bit}_j(x_{is}+512)w_{10s+j}
       +\sum_{j=40}^{43}w_j .
\]
Here \(\operatorname{bit}_j(m)=\lfloor m/2^j\rfloor\bmod2\)
is used only to define the finitely many fixed integer coefficients.
It is not an operation on a quantified Diophantine variable.
The other coefficients are
\[
c_{e_s(j)}=w_{10s+j},\quad
(c_{1095},c_{1096},c_{2195},c_{3296})=(w_{40},w_{41},w_{42},w_{43}),
\quad c_{1063}=c_{1064}=0.
\]
Define
\begin{align*}
\mu&=1022+\sum_{j=0}^{43}w_j,\\
Z_{\rm pad}&=(2^{119}-1)\sum_{e=8792}^{8798}\rho^e
                +(2^{27}-1)\rho^{8799},\\
D_C&=\rho^{40}+\rho^{41}+\rho^{1100}+\rho^{2200}
       +\sum_{e\in\mathcal E}c_e(\rho^{4393-e}+\rho^{5495-e}),\\
D_R&=\rho,\\
M_C&=B_0-1-\sum_{\substack{e\in\mathcal E\\e\ne1}}\rho^e
                   -2\rho^{1063},\\
M_{F,\rm native}
 &=4+\mu(\rho^{4393}+\rho^{5495})
     +\sum_{j=0}^9(\rho^{1000+2j}+\rho^{1021+2j})
     +\rho^{2195}+\rho^{3296}+Z_{\rm pad}.
\end{align*}
The six actual fixed ports are precisely
\[
\boxed{\begin{gathered}
\beta=B_0-1,\quad K_0=D_C+B_0D_R,\quad M_C,\quad
M_F=M_{F,\rm native}+B_0-1,\\
A_{\rm cone}=2^{1113},\qquad B_{\rm cone}=2^{666104}.
\end{gathered}}
\]
All displayed exponents are fixed integers. In particular,
\(M_F\) already includes its shift; shifting it again would change
the equation and invalidate the binding.

\begin{table}[ht]\centering
\begin{tabular}{lrr}\toprule
Fixed numeral & Bit length & Minimal unsigned bytes\\\midrule
\(\beta\)&1\,064\,921&133\,116\\
\(K_0\)&1\,065\,043&133\,131\\
\(M_C\)&1\,064\,921&133\,116\\
\(M_F\)&1\,064\,922&133\,116\\
\(A_{\rm cone}\)&1\,114&140\\
\(B_{\rm cone}\)&666\,105&83\,264\\\midrule
Total&&615\,883\\\bottomrule
\end{tabular}
\caption{Exact fixed payload sizes before compression.
The archive stores deterministic gzip encodings of the minimal
unsigned little-endian bytes and binds their raw hashes.}
\end{table}

The source proof verifies the whole coefficient interface, not just
these sizes. It establishes local tile decoding, carry-free geometry,
whole-cell rotation, mask populations, positive canonical slack, and
an explicit rank controller. In particular
\[
\operatorname{pc}(M_C)+\operatorname{pc}(M_{F,\rm native})=d,\quad
M_C\equiv2\pmod4,\quad M_{F,\rm native}\equiv4\pmod8.
\]
The controller uses
\(d=11^2\cdot13\cdot677\), \(\lambda=2341\), and a proved Boolean
subset coverage lemma; it does not assume that \(d\) is a power of
five. Its construction supplies odd canonical exponents, odd cubic
scale, and the canonical parity promises in \eqref{eq:canonical}.
The complete inherited proof and literal builder are included,
with their internal source pins and licenses preserved.



\begin{thebibliography}{99}
\bibitem{repo} V. Reshetnikov, \emph{ProveIt}, reviewed commit
\texttt{04a1d79222d55cabe4ba52897863fc33baa8c268}.
\url{https://github.com/VladimirReshetnikov/ProveIt/tree/04a1d79222d55cabe4ba52897863fc33baa8c268/Computability/DiophantineConsistency}.
\bibitem{fusion} \emph{Thirteen Integer Variables through a Single
Five-Radical Norm}, incoming package
\texttt{ConZFC\_13\_Integer\_Norms\_2026-10-24.zip}, in \cite{repo}.
Complete package retained under \texttt{inherited/root\_fusion}.
\bibitem{coeff} ProveIt source-34 finite coefficient recipe and ordinary
compiler proof, retained in \cite{fusion},
\path{provenance/upstream_ordinary/source34_coefficients.md} and
\texttt{coefficients/recipe/}.
\bibitem{signed} ProveIt ordinary and signed Pell/native proofs, retained
in \cite{fusion}, \path{provenance/upstream_ordinary/ordinary36.md},
\path{provenance/RESIDUE_AND_DETERMINANT.md}, and
\path{provenance/NORMALIZED_NORMS_PROOF.md}.
\bibitem{guards} ProveIt, \emph{Joint guards and signed residue
orientation}, incoming
\texttt{ZFC\_Diophantine\_Joint\_Guards\_2026-10-24.zip}.
The two relevant proofs and their notices are retained in
\path{provenance/signed_guards/}.
\bibitem{calibration} ProveIt ordinary K215 consistency calibration,
retained in \cite{fusion},
\path{provenance/upstream_ordinary/k215_ordinary.md}.
\bibitem{active} ProveIt, \emph{A Diophantine equation for syntactic
consistency of ZFC}, active research README at the reviewed commit.
\bibitem{boundary} ProveIt, \emph{Object-ZFC formalization status}, as
retained in the incoming package's
\path{provenance/CALIBRATION_STATUS.md} and pinned source ledger.
\bibitem{sun2026} Z.-W. Sun, \emph{Ten unknowns for Hilbert's tenth problem
over the integers}, arXiv:2609.01594v2, 2026.
\url{https://arxiv.org/abs/2609.01594}.
\bibitem{sanna} C. Sanna, \emph{The $p$-adic valuation of Lucas sequences},
Fibonacci Quarterly \textbf{54} (2016), 118--124. The valuation results are
also explicitly stated and applied in \cite{sun2026}, Lemmas 2.3--2.4.
\end{thebibliography}
\end{document}
